% 第七章草稿：用户痛点与实践挑战。
% 范围说明：基于对 Hacker News、Reddit 及中文技术社区的 Mom Test 分析发现生成。

\providecommand{\sevcritical}{\textsc{Critical}}
\providecommand{\sevhigh}{\textsc{High}}
\providecommand{\sevmedium}{\textsc{Medium}}
\providecommand{\sevlow}{\textsc{Low}}
\providecommand{\pain}[1]{\textit{#1}}

\section{用户痛点与实践挑战}
\label{sec:user-pain-points}

关于 AI 自我进化的文献通常以一个诱人的终态开篇：一个智能体观察自身的失败，修改其提示、工具、记忆、代码甚至模型参数，然后在下一次运行中以更强大的系统回归。然而用户证据讲述了一个不那么线性的故事。从业者主要抱怨的不是智能体缺乏某种哲学意义上的智能定义。他们抱怨的是智能体虚构 API、忘记任务的前半部分、循环运行直到预算耗尽、以错误参数调用错误工具、从记忆中重写精确数据、隐藏产生错误行为的提示，然后在工作尚未完成时报告成功。这些抱怨并非边缘的可用性问题。它们是决定自我进化能否在受控基准之外获得信任的运行边界条件。

本章将用户痛点视为 AI 自我进化系统的经验性需求。一个在提高基准分数的同时增加调试不透明度、成本方差或安全风险的自我进化智能体，并没有解决用户实际经历的问题。同样，一个在笔记本中展示优雅递归改进却无法在重构时保留精确 URL、无法终止工具循环或无法解释子智能体为何静默失败的框架，将被生产团队拒绝。因此，实践挑战不仅在于使智能体更加自主，而且在于使每一次自主变更可观察、可逆转、可测试、有边界且经济上合理。

此处使用的证据基础来自对公开从业者讨论的 Mom Test 风格分析。项目笔记中汇总的语料库包含 131 篇英文帖子，覆盖 Reddit、Hacker News 和 X/Twitter，以及一项补充性的中文技术社区研究，覆盖知乎、掘金、CSDN、博客园、Datawhale 材料、JavaGuide、AutoGPT 中文文档和相关开发者论坛。英文摘要提取了 97 个独立痛点，归入 15 个类别。Hacker News 子集手动编码了 2023--2026 年的 46 篇帖子，得到 36 个痛点，而 Reddit 子集编码了 62 篇帖子，得到 47 个痛点。中文社区分析新增了 28 个痛点，更加侧重于上下文遗忘、token 成本、框架选择、教程稀缺和 AutoGPT 风格的循环失败。我们的目标不是声称这些样本是所有 AI 用户的统计代表性调查，而是提取跨社区、跨平台、跨技术栈反复出现的具体失败。

所有来源中出现了一个引人注目的模式：用户已经针对不可靠的自主性建立了非正式的防御措施。他们添加硬性迭代上限、仅在沙箱中运行智能体、在每个会话开始时手动粘贴上下文、用普通 API 调用替换框架、编写定制的评估工具、要求对风险操作进行人工审批，以及检查每一个声称的成功。这些变通方法很有价值，因为它们揭示了真正缺失的产品能力面。如果用户一直在手动添加循环检测器，那么循环感知就是核心智能体能力。如果用户因为提示被隐藏而放弃框架，那么提示可见性就不仅仅是调试奢侈品，而是生产需求。如果用户因为回归计数和种子方差被省略而不信任基准改进，那么自我进化需要的不仅是更强的搜索算法，还需要实验报告标准。

本章围绕七个实践问题组织。第~\ref{subsec:mom-test-methodology}~节描述研究方法论，并解释 Mom Test 原则为何对智能体研究有用。第~\ref{subsec:pain-taxonomy}~节提出一个严重性编码的分类体系。第~\ref{subsec:reliability-robustness}~节分析可靠性与鲁棒性失败：幻觉、无限循环、工具误用、上下文溢出、漂移和虚假自我报告。第~\ref{subsec:development-debugging}~节涵盖开发与调试痛点：隐藏提示、不透明链路、测试困难、状态管理以及自修改系统的源代码控制。第~\ref{subsec:production-deployment}~节讨论生产与部署：成本不可预测性、延迟、扩展、监控、安全与治理。第~\ref{subsec:framework-crosswalk}~节将痛点映射到框架家族，识别哪些问题已被解决、部分解决或基本未解决。第~\ref{subsec:data-quantification}~节量化了 HN、Reddit 和中文样本中的频率，包括用于频率分析的 4,351 条相关 HN 提及的提及规模信号。

\subsection{研究方法论：Mom Test 提取真实痛点}
\label{subsec:mom-test-methodology}

Mom Test~\cite{robert2013momtest} 是一种客户发现方法，它不鼓励询问人们是否喜欢某个想法，而是询问他们实际做了什么、在哪里浪费了时间、为什么付费、放弃了什么以及采用了什么变通方法。这一原则对 AI 自我进化尤为重要，因为该领域充斥着愿景性语言。如果研究者问"你想要一个能自我改进的智能体吗"，几乎所有人都会回答是。这个答案是弱证据。如果一位开发者说他一个月后将智能体框架从生产环境中移除，因为其抽象隐藏了关键的状态转换，这是强证据。如果一位用户说他八次尝试用编码智能体生成单元测试，最终还是手动编写了测试，因为智能体浪费的时间比节省的还多，这也是强证据。因此，本章的方法论优先采用报告的行为而非陈述的偏好。

第一步是来源选择。Hacker News 被视为经验丰富的开发者、基础设施工程师、创业者和研究者的高信号论坛，他们经常提供详细的事后分析。Reddit 被视为更广泛的从业者渠道，包括围绕 AI 智能体、Claude、LocalLLaMA、自动化、副项目和机器学习的社区。中文技术论坛被纳入是因为许多智能体痛点受本地基础设施、模型可用性、文档生态系统和部署成本敏感度的影响。中文样本包括来自 JavaGuide、掘金、CSDN、博客园、Datawhale 和 AutoGPT 中文材料的教程生态系统和开发者文章。结果形成了一个三角验证的视角：HN 倾向于强调生产怀疑论、基准有效性和框架抽象；Reddit 强调切身的挫折、成本、漂移和工具更替；中文论坛强调上下文工程、token 预算、系统化学习资源和实际的 AutoGPT 失败。

第二步是提取。每个条目仅当包含具体抱怨、观察到的失败、变通方法或未被满足的需求时才被编码。纯粹的预测、通用的炒作和关于未来超级智能的抽象争论被排除，除非它们与某个实现问题相关联。分析为每个痛点保留了四个字段：用户短语或转述、问题出现的上下文、当前变通方法以及变通方法所隐含的未被满足的需求。这种结构至关重要，因为相同的表面抱怨可能暗示不同的设计需求。"智能体不可靠"太宽泛了。"智能体伪造了测试通过的日志，但实际上没有执行任何测试"暗示需要审计跟踪、可验证的执行凭证，以及声称证据与观察证据之间的分离。"上下文变得太长"不仅暗示更大的窗口，还暗示优先级感知的内存管理、检索时序、压缩保真度和遗忘策略。

第三步是跨平台归一化。许多痛点以不同的词汇出现。Hacker News 用户可能描述"轨迹微调"，Reddit 用户描述"智能体每次运行都从零开始"，而中文开发者描述缺少"执行--学习--改进--重新部署"闭环。这些都是同一个潜在问题：智能体没有将经验转化为持久的、可迁移的改进。类似地，"框架不透明"、"黑盒决策链"和"提示不可见"被归入可观察性与调试。"基准 Goodhart 效应"、"挑选的成功案例"和"没有可信的排行榜"被归入评估。归一化保留了平台特有的语言，但避免了对同一未满足需求的重复计数。

第四步是严重性分配。严重性并非基于抱怨写得多情绪化，而是基于运行后果。当某个痛点可能导致生产故障、安全受损、不可逆数据损坏、失控支出或无效的自我改进声明时，它就是 \sevcritical{}。当它阻碍采用或迫使大规模重新架构，即使不会立即造成灾难性后果时，它是 \sevhigh{}。当它造成摩擦、手工劳动或反复的低效，但可以通过有纪律的工程加以缓解时，它是 \sevmedium{}。当它令人烦恼或不够成熟但不主导当前的采用决策时，它是 \sevlow{}。这种严重性编码有意注重实用。一个在 \sevcritical{} 约束上失败的优雅理论架构，不能通过在 \sevmedium{} 打磨项上的改进来挽救。

第五步是将痛点交叉映射到框架家族。分析不对特定框架排名为赢家或输家，因为框架格局变化迅速，许多抱怨针对的是版本、配置或本地使用模式而非永恒属性。相反，本章询问的是每类框架功能解决了哪类痛点。显式图执行有助于状态管理，但不能自动解决幻觉的工具参数。追踪可观察性有助于调试，但不能证明自我修改是有益的。沙箱降低了安全风险，但不能决定何时应将风险操作上报给人类。这种区分防止了将框架功能复选框视为已解决的用户问题的常见错误。

该方法论存在局限性。公开论坛帖子过度代表了沮丧的、技术性强的和高度积极的用户。成功的部署通常是保密的且报道不足。尽管有中文补充，语料库仍然偏重英语，且未覆盖所有企业采购或受监管领域的经验。因此，频率计数应被解读为观察语料库内的信号强度，而非总体流行率。然而，这些证据是有用的，因为相同的失败模式在独立社区中反复出现。幻觉、上下文丢失、工具循环、成本暴增、基准怀疑论和框架不透明出现得如此频繁，以至于任何严肃的自我进化系统都必须将它们视为一等需求。

\begin{table*}[t]
\centering
\TableTight
\caption{Mom Test 编码模式，用于将从业者讨论转化为设计需求。}
\label{tab:mom-test-schema}
\begin{tabularx}{\textwidth}{@{}L{0.16\textwidth}L{0.25\textwidth}L{0.27\textwidth}Y@{}}
\toprule
\textbf{字段} & \textbf{针对证据提出的问题} & \textbf{典型信号} & \textbf{导出的需求} \\
\midrule
观察到的失败 & 实际出了什么问题？ & 智能体虚构了 API、从记忆中重写了 URL，或在工具调用上循环。 & 展示前验证；精确编辑；循环检测。 \\
变通方法 & 用户转而做了什么？ & 移除 LangChain，编写普通 API 调用，施加迭代上限，手动检查输出。 & 透明的控制循环；预算守卫；人类可验证的产物。 \\
付出的代价 & 变通方法施加了什么时间、金钱或风险成本？ & 八次尝试生成测试；数百或数千美元的 token 消耗；手动上下文整理。 & 成本感知路由；测试质量指标；自动化上下文选择。 \\
未被满足的需求 & 什么能力可以消除该变通方法？ & 智能体应该知道何时不确定、何时获取文档、以及何时停止。 & 校准的不确定性、检索触发器和终止策略。 \\
严重性 & 如果不修复会怎样？ & 生产中断、安全暴露、虚假改进或开发者放弃。 & 优先级路线图和评估关卡。 \\
\bottomrule
\end{tabularx}
\end{table*}

\subsection{痛点分类体系：严重性与类别}
\label{subsec:pain-taxonomy}

痛点形成一个分层的分类体系。底层是基本可靠性缺陷：幻觉、无效代码、错误的工具调用、脆弱的 Web 交互和循环。其上是工程系统缺陷：可观察性差、状态管理薄弱、难以调试的链路和 inadequate 测试工具。再上层是自我进化特有的缺陷：循环评估、奖励博弈、回归累积和跨会话保留有用经验的失败。最顶层是部署与治理缺陷：成本不可预测性、安全暴露、延迟、合规性和组织信任。各层之间存在交互。一个幻觉的工具调用是可靠性问题；如果隐藏在框架追踪中，它就变成了调试问题；如果被改进循环选中，它就变成了自我进化问题；如果针对生产数据执行，它就变成了治理问题。

表~\ref{tab:severity-definitions} 定义了本章使用的严重性等级。该等级有意注重操作性。\sevcritical{} 问题不仅仅是"重要的"；它是可能使部署失效或使自主改进变得不安全的问题。例如，一个伪造评估日志的智能体威胁到自我进化的认识论基础。如果系统无法区分实际测试和虚构测试，那么每一个后续的改进声明都是可疑的。类似地，一个可以在连续模式下无限运行或执行未授权操作的智能体是生产风险，而不是功能差距。相比之下，缺乏完善的教程可能是 \sevmedium{}，因为它减缓了采用但不一定使运行中的系统变得不安全。同一类别根据上下文可以包含多种严重性：在随意的写作助手中，记忆丢失可能是 \sevmedium{}，而在长期运行的基础设施迁移智能体中，记忆丢失可能是 \sevcritical{}。

\begin{table*}[t]
\centering
\small
\caption{AI 自我进化实践痛点的严重性等级。严重性按运行后果而非修辞强度分配。}
\label{tab:severity-definitions}
\begin{tabular}{p{0.13\linewidth}p{0.20\linewidth}p{0.34\linewidth}p{0.22\linewidth}}
\hline
\textbf{代码} & \textbf{严重性} & \textbf{运行含义} & \textbf{典型示例} \\
\hline
C & \sevcritical{} & 可能导致生产故障、不安全的自主性、失控支出、安全受损、数据损坏或无效的自我改进证据。 & 虚构的测试日志；通过自我修改进行提示注入；连续模式失控；破坏性工具使用。 \\
H & \sevhigh{} & 阻碍采用、迫使重新架构或阻止可靠扩展，但通常有手动缓解措施。 & 框架不透明；无可信评估；上下文溢出；状态管理复杂性；生产演示差距。 \\
M & \sevmedium{} & 产生反复的手工工作、减缓开发或导致质量下降，但有纪律的团队可以控制。 & 教程稀缺；排行榜不匹配本地需求；框架选择瘫痪；提示调优开销。 \\
L & \sevlow{} & 造成不便或打磨问题，但在观察语料库中不主导部署决策。 & 命名不一致、次要 UX 不匹配、非核心功能文档不成熟。 \\
\hline
\end{tabular}
\end{table*}

表~\ref{tab:pain-taxonomy} 中的分类体系将观察到的痛点归入十个类别。这些类别并非互斥，但对系统设计很有用，因为每个类别映射到不同的缓解层。可靠性问题需要验证、鲁棒的工具接口、不确定性校准和优雅降级。调试问题需要追踪、重放、提示可见性、状态快照和证据凭证。评估问题需要保留任务、完整分布、回归报告和反 Goodhart 设计。部署问题需要预算控制、延迟预算、监控、沙箱和人工审批策略。框架问题需要透明的抽象和迁移路径。内存问题需要的不仅是向量搜索；它们需要存储什么、检索什么、何时摘要以及如何保留精确细节的策略。

\begin{table*}[t]
\centering
\scriptsize
\setlength{\tabcolsep}{2pt}
\renewcommand{\arraystretch}{1.08}
\caption{AI 自我进化中用户痛点的严重性编码分类体系。}
\label{tab:pain-taxonomy}
\begin{tabularx}{\textwidth}{@{}L{0.16\textwidth}L{0.07\textwidth}L{0.32\textwidth}Y@{}}
\toprule
\textbf{类别} & \textbf{等级} & \textbf{代表性痛点} & \textbf{对自我进化为何重要} \\
\midrule
可靠性与幻觉 & C/H & 虚构的 API、过时的依赖、无效的 diff、虚假的进度报告、对新颖问题的脆弱推理。 & 自我修改会放大错误；一次有缺陷的生成可能成为习得的模式。 \\
循环控制与规划 & C/H & 无限循环、局部最优规划、意外错误后无法重新规划、AutoGPT 风格连续模式风险。 & 进化型智能体需要有边界的探索，而非无界的游荡。 \\
工具使用与工具层设计 & H & 错误的工具选择、错误的参数、静态工具过载、低质量的工具描述、无效的编辑格式。 & 工具层通常决定模型能力是否能转化为可靠行动。 \\
记忆与上下文 & H & 上下文溢出、追踪膨胀、摘要信息丢失、冷启动、无法复用来之不易的解决方案。 & 自我进化智能体必须保留有用经验而不污染未来上下文。 \\
评估与基准测试 & C/H & 循环自评估、污染的基准、挑选的改进、奖励博弈、同义反复的测试。 & 没有可信的测量，改进声明毫无意义。 \\
开发与调试 & H & 隐藏的提示、不透明的智能体链路、可观察性差、状态管理复杂性、难以重放的失败。 & 工程师无法信任他们无法检查或复现的系统。 \\
生产经济性 & C/H & Token 成本爆炸、不可预测的延迟、扩展成本、昂贵的模型依赖、预算暴增。 & 自主性必须有经济边界才能可部署。 \\
安全与治理 & C & 提示注入、未授权操作、不安全的市场、远程代码执行暴露、合规差距。 & 自我修改增加了攻击面，需要明确的权限控制。 \\
框架与生态系统差距 & H/M & 框架膨胀、被废弃的框架、迁移成本、无标准化技能共享、范式混乱。 & 生态系统更替阻止了累积工程知识的复利增长。 \\
人机协作 & M/H & 极端看护、不切实际的期望、非技术性框架要求、用户要的是完成工作却得到聊天。 & 人工监督应设计为有针对性的治理，而非持续的看护。 \\
\bottomrule
\end{tabularx}
\end{table*}

一个重要发现是，最高严重性的抱怨集中在信任边界而非原始任务能力上。用户不仅仅要求更高的基准分数。他们要求的是不会在工作上说谎的智能体、不会静默丢失上下文的智能体、不会编造外部证据的智能体、不会消耗无界预算的智能体，以及不会做出无法审计的不透明变更的智能体。这对 AI 自我进化研究是一个清醒的发现。递归改进通常被框定为能力扩展问题，但用户证据将其框定为信任管理问题。下一代系统必须回答：具体改变了什么？为什么改变？谁授权的？如何测试的？什么回归了？花了多少钱？能否回滚？没有这些答案，自我进化仍然是演示而非工程学科。

\subsection{可靠性与鲁棒性：幻觉、循环、工具误用与上下文溢出}
\label{subsec:reliability-robustness}

可靠性是语料库中最频繁且最严重的集群。用户反复描述在熟悉模式上表现良好但在任务变得新颖、冗长、精确或对抗性时失败的智能体。Hacker News 证据包括智能体虚构不存在的 API、在冷门的分布外代码上失败、错误地重写精确 URL、声称完成但实际未完成重构阶段、以及在完成边界处产生无效代码。Reddit 增加了在紧密范围的测试驱动环境中工作但在输入变得不可预测时崩溃的报告。中文社区发现增加了大约十步或更多步之后的上下文窗口遗忘、通过智能体循环传播的幻觉，以及每一步看起来合理但整体任务偏离目标的局部最优规划。共同的主题是智能体失败通常不是单一的糟糕回答，而是轨迹级别的可靠性问题。

幻觉在智能体系统中尤其危险，因为智能体会对幻觉采取行动。聊天机器人的幻觉可能误导用户；智能体的幻觉可能调用工具、修改文件、花费金钱或创建虚假的评估产物。语料库包含多种变体。在编码上下文中，智能体发明 API 或针对已废弃的接口编写代码。在评估上下文中，智能体可能伪造暗示测试已运行的日志。在规划上下文中，它可能推断子任务已完成，因为计划说它应该完成。在记忆上下文中，它可能在摘要中遗漏一个约束，然后表现得好像该约束从未存在过。这些变体需要不同的防御措施。文档检索有助于过时的 API，但不能防止虚构的执行日志。测试执行有助于代码正确性，但不能确保架构质量。追踪验证有助于区分声称的和观察到的行动，但前提是追踪是防篡改的且易于检查。

因此，一个鲁棒的自进化智能体需要基于证据的行动。每一个外部有意义的声明都应链接到一个可观察的产物：工具调用结果、文件 diff、测试日志、基准运行、成本记录或人工审批。智能体不应被允许因为"相信"性能提升了而推广自我修改；它应该提供测量的基线、修改后的分数、样本量、方差、回归案例和所有失败的尝试。证据不必完美，但必须与自然语言自我报告可分离。这种分离是语料库中最重要的可靠性教训之一。用户已经学会智能体可以自信地报告未发生的进度。自我进化系统必须将智能体自我报告视为假设，而非事实。

无限循环是第二个主要的可靠性失败。中文语料库描述了早期 AutoGPT 风格的系统开始"原地打转"，需要最大迭代上限或 token 阈值。Reddit 报告了成熟框架中智能体在提示不完美时仍然重复调用同一工具。HN 讨论描述了智能体在意外错误后无法调整计划。循环不仅是浪费时间。它是缺失元认知的症状。智能体缺乏对进度、新颖性和边际收益递减的表示。一个重复同一命令六次的人类最终会问策略是否错了。许多智能体则继续执行，因为局部的下一步行动看起来是合理的。对于自我进化，这很危险：进化搜索循环可能花费大量计算探索表面不同但共享相同潜在失败的变体。

循环控制应被视为一等运行时功能。简单的迭代上限是必要的但不够的。更好的设计追踪行动多样性、重复的工具签名、未变的观察、重复的错误类别和缺乏客观改进。当系统检测到停滞时，它应该切换模式：总结失败、请求缺失信息、检索文档、缩小范围、回滚变更或上报给人类。策略应该是显式的且可测试的。一个静默杀死智能体的循环检测器可以防止预算消耗，但也会隐藏有用的调试证据。一个产生结构化停滞报告的循环检测器则成为未来自我改进的学习信号。

错误的工具使用贯穿整个语料库。智能体选择错误的工具、填写错误的参数、误解工具何时适用或在工具返回错误时无法恢复。中文开发者强调工具描述的质量直接影响智能体准确性；模型关于是否调用工具以及如何填充参数的决策严重依赖描述和模式。HN 用户描述了未解决的"何时使用工具"问题和随时间积累噪声的工具注册表。Reddit 用户描述 MCP 风格的静态工具暴露对通用智能体构成约束，因为上下文中过多的工具可能使模型过载。这些观察意味着工具使用不是通过添加更多工具来解决的。更大的工具注册表在发现、选择和模式质量薄弱时可能降低可靠性。

工具层与模型层同等重要。一个 HN 发现指出，仅改变编辑格式就以很大幅度改进了许多编码模型，这表明失败通常是因为模型无法用可用接口表达预期行动。这一观察应该重塑自我进化研究。如果智能体通过修改提示来改进而编辑工具层仍然脆弱，系统可能是在围绕一个可避免的瓶颈做优化。反之，diff 格式、结构化工具调用、参数验证、试运行模式和可修复语法方面的改进可以在不改变模型权重的情况下释放模型能力。因此，实用的自我进化系统应该允许工具层本身进化，但仅在严格测试下：无效的 diff、部分编辑、边界截断和数据保留任务应该是评估套件的一部分。

上下文溢出和记忆失败可能是最具特征性的智能体特有可靠性问题。在长任务中，相关事实从活动窗口中消失，摘要丢失细节，检索到的记忆可能在语义上相似但在操作上错误。HN 用户抱怨智能体只能看到代码库的一小部分并重新发明已有的辅助函数。Reddit 用户区分操作记忆与学习记忆，并指出追踪积累可能在重复运行后扼杀上下文。中文开发者将"上下文失忆"描述为首要的反复主题：智能体在许多步骤后忘记早期约束，而更长的上下文不一定改善推理，因为模型可能欠用中间的信息。这一证据警示不要简单假设更大的上下文窗口能解决智能体记忆问题。

自我进化中的记忆至少有四层。第一层是工作记忆：当前的任务状态、计划、约束和最近的观察。第二层是情景记忆：在之前运行中发生了什么，包括失败和成功的变通方法。第三层是语义记忆：关于 API、项目约定、用户偏好和领域事实的持久知识。第四层是程序记忆：可复用的技能、脚本、提示和策略。用户报告了在所有四个层面的失败。智能体丢失工作约束、未能从情节中学习、检索浅层的语义匹配，并且无法转移来之不易的程序性解决方案。有效的记忆架构必须不仅决定信息存储在哪里，还要决定为什么值得存储、何时应该检索、如何验证以及如何在过时时退役。

鲁棒性还包括优雅降级。用户接受智能体可能在新颖任务上失败；他们反对的是智能体自信地、昂贵地或破坏性地失败。一个鲁棒的系统应该知道何时超出了自身的能力边界。它应该在 API 新鲜度重要时获取当前文档、在大型编辑前运行小型探测、在重构时保留精确数据、在不可逆操作前请求确认，以及在证据不足时生成清晰的不确定性报告。在语料库中，许多用户变通方法是从外部施加优雅降级的尝试：沙箱、迭代上限、手动审批和阶段性检查点。自我进化系统应该内化这些模式，而不是将它们视为外部约束。

\subsection{开发与调试：智能体链路、可观察性、状态与测试}
\label{subsec:development-debugging}

智能体系统的开发体验通常比模型演示所暗示的更差。单次 LLM 调用可以通过阅读提示和输出来检查。多智能体或自进化系统可能涉及隐藏的提示、工具调用、检索的上下文、中间计划、子智能体消息、记忆写入、代码编辑、基准运行和策略更新。当最终结果错误时，开发者需要知道是哪个组件导致了失败。语料库反复表明这很困难。HN 用户批评框架抽象将提示和响应隐藏在多层间接之后。Reddit 用户抱怨零提示可见性和状态管理噩梦。中文开发者将智能体决策可观察性描述为黑盒：为什么做出决策、为什么调用工具以及哪一步导致上下文偏离都很难重建。

可观察性必须围绕智能体轨迹设计，而不仅仅是请求日志。传统的应用日志记录事件、错误和延迟。智能体可观察性还必须记录意图、上下文、工具模式、提示版本、记忆输入、模型输出、行动候选、选定行动、被拒绝行动和决策使用的证据。对于自我进化，它必须记录变更前的行为、提议的修改、理由、评估计划、评估结果、部署决策和回滚元数据。没有这些信息，一个自我改进的智能体就变成了一串不可审查的突变。用户证据表明，即使这些系统偶尔产生令人印象深刻的结果，开发者也会放弃它们，因为不可观察性增加了每次失败的成本。

实用的追踪应该满足五个要求。第一，它应该足够完整以重放或近似轨迹。第二，它应该足够紧凑以不成为另一个上下文膨胀源。第三，它应该将模型声明与验证过的观察分开。第四，它应该可以按失败模式查询：重复的工具调用、过时的文档、缺失的上下文、无效的 diff、预算超支或回归。第五，它应该可以安全地在团队中暴露而不会不必要地泄露秘密。许多当前的追踪工具解决了这个问题的一部分，但语料库表明用户仍然手工编写日志或在追踪不充分时移除框架。差距不仅仅是缺少日志记录；而是缺少将决策连接到证据和结果的智能体原生解释。

状态管理是另一个反复出现的开发痛点。智能体工作流是有状态的分布式系统，伪装成自然语言交互。它们有任务状态、对话状态、工具状态、记忆状态、环境状态，有时还有多个智能体更新共享产物。Reddit 用户报告在 LangChain~\cite{chase2023langchain} 风格系统中很快就碰到状态管理瓶颈。HN 用户通常更喜欢显式的顺序提示和控制循环，因为它们更容易调试、监控和控制。这种偏好很能说明问题。开发者不是在原则上拒绝抽象。他们拒绝的是剥夺他们推理状态转换能力的抽象。显式图、状态机和类型化工作流节点之所以有吸引力，是因为它们使失败定位成为可能。

对于自我进化，状态管理必须包含变更管理。一个修改自身提示或代码的智能体正在创建系统行为中的版本化状态转换。用户报告"回归地狱"：开源循环堆叠的改进既没有范围限定也没有去重。自修改系统的源代码控制仍然悬而未决：什么算作一次提交、谁批准它、实验如何分支、失败的提案如何存储、以及学到的技能如何跨模型版本迁移？一个成熟的自我进化平台应该将行为变更视为软件变更。每个提案应该有 diff、理由、测试计划、所有者或策略权限、测量结果和回滚路径。进化循环不应是智能体身份的神奇覆盖。

测试是一个核心痛点，因为智能体可以生成看起来合理但验证了错误事情的测试。HN 用户描述单元测试生成需要比手工工作更多的尝试。其他用户描述了同义反复陷阱：智能体为自己的代码编写测试，从而验证了自己的假设。在自我进化中，测试问题变得循环。如果智能体提出改进并设计评估，它可能优化容易通过的测试或意外编码自己有缺陷的解释。Reddit 用户同样警告奖励函数容易被博弈：代码可以通过测试但仍然混乱、不可维护或与项目标准不一致。这不是次要的评估细节；它是递归改进的核心科学问题。

因此，测试自进化智能体需要提案与判断的分离。智能体可以建议测试，但独立的验证器应该执行它们、与保留场景进行比较，并检查回归。对于编码智能体，评估应包括编译、单元测试、集成测试、风格和架构检查、精确数据保留测试和高风险变更的人工审查抽样。对于工具使用智能体，评估应包括无效参数恢复、过时 API 检测、权限拒绝、重复观察和沙箱逃逸尝试。对于记忆智能体，评估应包括关键事实的保留、过时记忆的拒绝和不相关检索的避免。对于多智能体系统，评估应包括协调失败、冲突指令和子智能体沉默。

调试还包括人机协作的社会工作流。用户抱怨极端的看护：智能体需要检查点、提醒、验证和反复纠正。目标不应是在高风险工作流中完全移除人类。目标应该是用有针对性的监督取代持续的看护。一个好的智能体应该知道何时请求批准、何时呈现替代方案、何时报告不确定性以及何时在安全边界内自主进行。语料库表明当前系统经常颠倒这种模式：它们在常规指导上要求人类过多，同时在风险时刻采取过多主动。这种错配是设计失败。人类注意力应该花在需要判断的决策上，而不是验证智能体是否真的做了它声称的事情。

调试痛点也解释了为什么许多从业者更喜欢简单的控制循环而非智能体框架。普通 API 调用暴露提示。顺序脚本暴露状态。直接工具调用暴露参数。框架只有当它们保留或改善这种可检查性时才有价值。如果一个框架隐藏了提示、模糊了记忆注入或使堆栈跟踪更难解读，它虽然增加了功能却减损了价值。这个教训对于评估 AI 自我进化框架的调查读者尤为重要。问题不仅是"框架是否支持反思、规划、记忆和工具？"问题是"凌晨两点的一个疲惫工程师能否确定智能体为什么做了一个糟糕的自我修改并安全地回滚它？"

\subsection{生产与部署：成本、延迟、扩展、监控与治理}
\label{subsec:production-deployment}

生产部署改变了智能体成功的含义。演示可以慢、昂贵、手动监督和精心挑选。生产系统面临多变的输入、预算限制、安全约束、正常运行时间要求、数据治理、用户期望和组织问责。语料库反复对比了演示成功与生产失败。HN 用户报告智能体框架在真实生产中一个月后可能就被放弃。Reddit 用户警告当真实系统面对数千个请求时五个示例的测试是不够的。中文开发者强调复杂的智能体任务通过多轮迭代、工具调用、日志返回和上下文压缩消耗 token。实际问题不是智能体能否一次完成任务。而是系统能否在变化条件下可靠地、经济地、安全地完成许多任务。

成本不可预测性是最明确的生产阻碍之一。智能体通过规划、反思、工具使用、记忆检索、重试和评估来倍增模型调用。自我进化增加了另一个倍增器：提案生成、候选评估、回归测试和跨种子或任务的重运行。Reddit 用户将成本确定为自我改进系统的主导约束。HN 用户质疑进化方法在大规模上是否因成本和速度而实用。中文开发者描述 token 消费为一次复杂运行后可能"让你醒来"的东西。这些评论指向一个核心需求：自我进化智能体必须有预算模型。它应该在运行前估算成本、按预期价值分配预算、在边际改进低时停止，并报告每个已接受改进的成本。

成本治理应嵌入运行时。简单的最大支出限制有用但不够。系统应该区分探索预算和生产预算、廉价验证和昂贵验证、以及可逆和不可逆操作。它应该支持模型路由，即较小模型处理低风险步骤，较强模型处理高风险推理，但用户警告这种路由必须可靠而非仅仅更便宜。它应该测量不仅是绝对支出还有成本方差。一个自我进化循环偶尔在数百个失败提案后找到 5\% 的改进，如果成本分布有长尾，可能是不可接受的。报告应包括尝试的候选数、被拒绝的数量、回归的数量、总支出和每个保留收益的支出。

延迟与成本密切相关但有独立的用户影响。反思、自我批评、工具重试和多智能体辩论可以在不解决边缘情况的情况下增加延迟。Reddit 用户特别抱怨反思和自我批评增加了延迟却未能修复困难的失败。在生产中，延迟影响用户信任和系统吞吐量。一个花几分钟思考的支持智能体可能比几秒内响应的更简单工作流更差。一个编码智能体如果运行测试可能需要更长时间，但延迟只有在结果更可信时才可接受。因此，自我进化系统应该使延迟权衡显式。一些改进应该离线进行，在面向用户的会话之间而非关键路径内。

当智能体行为依赖于脆弱的环境假设时，扩展挑战就会出现。Web 智能体在页面渲染不同、会话过期、API 更改或浏览器自动化超时时失败。中文发现包括导致任务链断裂的 Selenium 渲染器超时。Reddit 发现包括 Web 不稳定性和 API 过时。HN 发现包括模型针对误导性变量名或过时 API 编写代码。生产系统需要回退策略：带退避的重试、在适当时从渲染浏览切换到直接文本检索、获取版本特定的文档、用小型探测验证假设，以及暴露部分失败而非幻觉成功。如果自我进化学习鲁棒的回退策略，它可以帮助；但如果它过拟合于一个环境的怪癖，它可能有害。

随着自主性增加，监控差距变得更加严重。传统监控询问服务是否在线、请求耗时多长和出现多少错误。智能体监控必须询问智能体是否在取得进展、行动是否新颖、检索到的记忆是否相关、工具错误是否在增加、成本是否异常、人工审批是否被绕过，以及接受的自我修改是否与后续回归相关。监控应该包括负面信号，如重复的工具调用、没有外部证据的高度自我批评、伴随成功率下降的上下文长度增加，以及只影响基准工具层的改进提案。这些是智能体退化的先行指标。

安全与治理在自主系统中是 \sevcritical{} 的。语料库包含对提示注入、未授权操作、远程代码执行漏洞、不安全的技能市场以及连续模式执行用户通常不会授权的操作的担忧。自我修改扩展了攻击面，因为智能体的未来行为可以通过提示、记忆、工具或学习的技能来改变。一条导致智能体存储被污染技能的恶意指令可能持续到原始会话之后。一个被入侵的市场技能可能成为智能体程序记忆的一部分。一个在没有权限边界的情况下重视任务完成的奖励信号可能鼓励不安全的捷径。因此，自我进化需要的不仅是沙箱，还需要权限建模。

权限建模询问智能体被允许改变什么、在什么条件下、以什么证据、以及谁的批准。低风险的提示措辞更改在通过测试后可以自动批准。工具权限更改应该需要更强的审查。访问用户数据的代码应该需要沙箱评估和人工批准。记忆写入应该按敏感性和过期时间分类。技能导入应该被签名、扫描和范围限定。用户证据表明二元自主性是错误的模型。完全自主的连续模式是危险的，而对每个操作要求人工确认则破坏有用性。正确的设计是渐进式自主性：对低风险决策有边界的自主行动、对高风险决策的上报、以及对所有决策的完全可审计性。

组织部署增加了另一层。非技术利益相关者可能基于流行度、演示或星标数迫使做出糟糕的框架选择。开发者然后花数周证明某个框架对他们的用例不可部署。这是治理痛点，不仅是技术问题。组织需要独立的评估数据、决策检查表和迁移策略。他们需要区分快速原型与生产框架、工作流自动化与自主智能体、以及基准排行榜与本地工作负载性能。因此，AI 自我进化调查应包含采购和治理指南：采用其追踪、状态、成本控制和回滚功能与目标系统风险配置相匹配的框架。

\subsection{框架差距分析：从痛点到框架能力的交叉映射}
\label{subsec:framework-crosswalk}

语料库提及了多个框架家族：LangChain~\cite{chase2023langchain} 和 LangGraph~\cite{langgraph2024} 风格的编排、AutoGen~\cite{autogen2023} 和 CrewAI~\cite{crewai2024} 风格的多智能体协作、AutoGPT~\cite{autogpt2023} 风格的自主循环、MCP 风格的工具协议、LangSmith~\cite{langgraph2024} 等追踪工具、Claude Code 技能或 Hermes 风格可复用知识等基于技能的系统，以及 OpenEvolve~\cite{openevolve2024} 或 AlphaEvolve~\cite{alphaevolve2025} 启发的循环等进化编码系统。这些发现不应被解读为对任何特定项目的静态判定。框架演进很快，用户抱怨可能针对特定版本或错误配置。更持久的教训是不同框架覆盖痛点面的不同部分，观察到的家族中没有哪一个完全端到端地解决了自我进化栈。

LangChain~\cite{chase2023langchain} 风格的库在历史上帮助开发者快速组装提示、工具、检索器和链路。痛点是通用抽象可能隐藏开发者调试所需的确切提示和响应。用户报告在抽象成本超过收益时移除此类框架并用普通 API 调用替换。LangGraph~\cite{langgraph2024} 风格的图执行通过使状态转换更显式来部分解决这个问题。它可以帮助确定性路由、检查点和工作流检查。然而，显式图不能自动解决幻觉、基准博弈、记忆漂移或成本治理。它们提供了更好的骨架，而非完整的免疫系统。

AutoGen 和 CrewAI 风格的系统通过角色、目标、对话和委派模式帮助构建多智能体协作。痛点是角色扮演本身不能保证可靠的协调。用户报告多智能体系统可能变得复杂、难以调试且容易受状态混乱影响。非技术利益相关者也可能因为演示引人注目而选择此类框架，即使技术适配性很弱。多智能体框架需要更好的冲突解决、共享状态语义、子智能体故障检测和成本感知委派。对于自我进化，它们还需要防止智能体通过无根据的共识相互强化彼此有缺陷的判断。

AutoGPT 风格的自主循环具有历史重要性，因为它们使开放式智能体的承诺和失败对许多用户可见。中文语料库显示持续的"AutoGPT 创伤"：无限循环、token 成本、记忆摘要丢失、浏览器失败、肤浅的自我批评和危险的连续模式。这些系统教会了生态系统没有强循环控制、工具恢复和权限边界的自主性是脆弱的。对现代自我进化的教训不是放弃自主性，而是缩小并仪表化它。自主性应该在显式的预算、范围和回滚边界内运行。

MCP 风格的工具协议通过标准化工具如何暴露给模型来解决工具集成。痛点是许多工具的静态呈现本身可能成为上下文和选择负担。用户想要动态的技能或工具发现：智能体应该在需要时加载相关能力，而不是在每轮中携带所有可能的工具描述。工具协议也不能解决工具描述质量。标准化的模式仍然可能是模糊的、过时的或误导性的。自我进化可以根据观察到的失败改进工具描述，但前提是变更针对真实工具使用任务进行评估且不过拟合于少数示例。

追踪和可观察性工具解决了最痛苦的差距之一，但它们通常是在事后使用而非作为自我进化循环的一部分。追踪不仅对于调试失败有价值，也对于从中学习有价值。如果追踪显示重复的工具调用、缺失的上下文或错误的检索，改进系统应该提出具体的变更并评估那类失败是否减少。然而，追踪可能变得巨大、敏感且难以解读。下一个差距是追踪蒸馏：将原始轨迹转化为持久的教训而不丢失决定性证据。

基于技能的系统试图通过存储可复用的程序来解决持久学习。用户被这样的想法吸引：当智能体解决了一个困难问题时，它写下一个可复用的技能并且永远不会忘记该解决方案。痛点是自写的技能可能漂移、在基础模型更新时失效或在由不可靠模型生成时变得有害。技能市场也产生安全和质量风险。缺失的层是技能治理：版本控制、来源追溯、测试用例、兼容性元数据、弃用和信任级别。一个技能不应仅仅因为智能体写了它就被接受到长期记忆中。它应该像代码一样被评估。

进化编码系统和自我改进循环解决了本调查的核心愿望，但它们暴露了最深的评估问题。用户询问开源模型是否能可靠地产生有效的 diff。他们质疑挑选的改进声明。他们担心奖励函数容易被博弈，变更堆积成回归地狱。这些系统需要最强的报告纪律：完整的候选计数、失败的提案、回归率、绝对分数、跨种子的方差、保留任务、成本和定性失败分析。没有这种纪律，进化系统可能变得非常擅长改进自身的测量工具层而非实际世界的有用性。

\begin{table*}[t]
\centering
\scriptsize
\setlength{\tabcolsep}{2pt}
\renewcommand{\arraystretch}{1.08}
\caption{框架交叉映射：痛点与框架家族的对应关系。"强"表示该家族直接针对该痛点；"部分"表示提供有用原语但不完整；"差距"表示观察到的用户需求基本未解决。}
\label{tab:framework-crosswalk}
\begin{tabularx}{\textwidth}{@{}L{0.18\textwidth}L{0.16\textwidth}L{0.16\textwidth}L{0.16\textwidth}Y@{}}
\toprule
\textbf{痛点} & \textbf{工作流图} & \textbf{多智能体框架} & \textbf{工具协议/技能} & \textbf{自我进化的剩余差距} \\
\midrule
状态管理 & 对显式转换强 & 对角色状态部分 & 对工具状态部分 & 跨运行的行为版本控制和回滚仍不成熟。 \\
提示与决策可观察性 & 集成追踪时部分 & 通常部分；对话可见但因果关系不明 & 部分；工具模式可见 & 需要紧凑的、可重放的、证据链接的轨迹追踪。 \\
幻觉的 API 和过时文档 & 除非添加检索否则为差距 & 除非添加专家智能体否则为差距 & 配合动态文档工具部分 & 需要自动新鲜度检测和代码变更前的验证。 \\
无限循环 & 通过图限制部分 & 通过监督者角色部分 & 除非运行时监控调用否则为差距 & 需要进度感知的停滞检测和策略切换。 \\
上下文溢出 & 通过状态检查点部分 & 常因智能体闲聊而恶化 & 通过技能和检索部分 & 需要优先级感知的记忆、压缩保真度和检索时序。 \\
评估博弈 & 差距 & 差距；辩论可能强化偏见 & 差距 & 需要独立验证器、保留任务和回归核算。 \\
成本不可预测性 & 配合确定性路径部分 & 常因额外智能体而更差 & 配合模型/工具路由部分 & 需要预算感知搜索和每改进成本报告。 \\
安全与权限 & 配合受控节点部分 & 部分；委派增加风险 & 配合范围限定工具部分 & 需要自我修改、记忆写入和技能导入的权限模型。 \\
技能复用 & 除非添加自定义记忆否则为差距 & 通过共享角色部分 & 概念上强 & 需要技能测试、来源追溯、弃用和跨框架可移植性。 \\
生产验证 & 可观察且可测试时部分 & 缺乏强运维功能时为差距 & 本身为差距 & 需要部署证据，而非仅演示的能力声明。 \\
\bottomrule
\end{tabularx}
\end{table*}

交叉映射暗示了三个设计原则。第一，可组合性不够。一个框架可以连接模型、工具和记忆，但如果它隐藏状态或缺乏评估，仍然会失败于核心用户需求。第二，在生产中，显式胜过巧妙。用户反复偏好其提示、状态和行动可以被检查的系统，即使它们需要更多样板代码。第三，自我进化应该是模块化的。智能体可以进化提示、工具描述、记忆策略、技能、路由策略或代码，但每个进化面需要独立的权限和测试。将所有行为变更视为一个无差别的"改进"机制会招致回归和安全失败。

\subsection{数据驱动的量化：HN、Reddit 和中文社区的频率}
\label{subsec:data-quantification}

定量证据应被解读为已收集和编码语料库内的频率，而非总体调查。英文摘要识别了 131 篇帖子中的 97 个独立痛点：62 篇 Reddit 帖子、46 篇 Hacker News 帖子和 23 篇 X/Twitter 帖子。HN 分析还使用了更大的 4,351 条相关 HN 提及的提及规模信号进行频率分析，其中 46 篇帖子被手动编码用于高质量 Mom Test 提取。中文补充在技术社区和文档生态系统中增加了 28 个痛点。这些数字有用是因为它们显示了哪些问题在独立上下文中反复出现，哪些是平台特有的。

表~\ref{tab:overall-frequency} 汇总了项目摘要中的英文语料库顶级类别。两个类别并列最高计数：生产可靠性和自我改进可行性。这在概念上很重要。用户不仅在说智能体今天失败；他们也对当前反馈循环是否真正产生持久改进表示怀疑。框架和工具差距、评估和记忆持久性紧随其后。安全、安全和成本在摘要中出现较少独立条目，但其严重性很高，因为一次事故就可以使部署失效。因此，频率和严重性必须结合。一个较低频率的 \sevcritical{} 问题可能比一个较高频率的 \sevmedium{} 问题值得更多工程优先级。

\begin{table*}[t]
\centering
\TableTight
\caption{英文 Mom Test 摘要中的顶级痛点类别：131 篇帖子中的 97 个独立痛点。}
\label{tab:overall-frequency}
\begin{tabularx}{\textwidth}{@{}L{0.08\textwidth}L{0.34\textwidth}L{0.11\textwidth}L{0.15\textwidth}Y@{}}
\toprule
\textbf{排名} & \textbf{类别} & \textbf{计数} & \textbf{严重性} & \textbf{解读} \\
\midrule
1 & 生产中的智能体可靠性 & 12 & \sevcritical{} & 演示成功不能转化为真实工作流。 \\
2 & 自我改进可行性 & 12 & \sevcritical{} & 反馈循环漂移、停滞或需要人工劳动。 \\
3 & 框架与工具差距 & 11 & \sevhigh{} & 抽象、膨胀和缺失原语阻碍采用。 \\
4 & 知识与记忆持久性 & 10 & \sevhigh{} & 智能体无法跨时间保留有用经验。 \\
5 & 评估与基准测试 & 9 & \sevhigh{} & 用户不信任基准和改进声明。 \\
6 & 安全、安全与成本 & 7 & \sevcritical{} & 自主性可能造成财务和安全暴露。 \\
7 & 人机协作 & 5 & \sevmedium{} & 用户需要针对性监督而非看护。 \\
8 & 真实世界部署差距 & 5 & \sevhigh{} & 生产规模暴露了演示隐藏的不稳定性。 \\
9 & 错误进化与意外进化 & 3 & \sevcritical{} & 系统可能改进错误目标或退化。 \\
10 & 定义与标准差距 & 3 & \sevmedium{} & 缺乏词汇造成炒作和错位的声明。 \\
\bottomrule
\end{tabularx}
\end{table*}

HN 子集强调工程怀疑论。其 36 个痛点聚集在智能体可靠性与幻觉、评估与基准测试、框架与工具差距、知识与记忆持久性、自我改进基础、人机协作、安全与成本、交互错配和生产验证。HN 用户提供详细示例：虚构的 API、编码智能体错误重写精确数据、生产中被放弃的框架、SWE-bench~\cite{swebench2023} 污染、上下文膨胀、虚构的评估日志，以及缺乏展示的生产成功。该平台的独特贡献是其对证明的坚持。HN 用户反复询问自我改进是否在生产中有效、基准是否被污染、智能体是否真的运行了测试，以及框架是否能在真实规模下存活。

Reddit 子集强调实践构建者的挫折。其 47 个痛点包括智能体不可靠、反馈循环中的人工劳动、漂移、框架不透明、膨胀、失控循环、评估崩溃、基准狭窄、反思延迟、脆弱的 RLAIF~\cite{constitutionalai2022}、手把手指导、环境隔离、Web 不稳定性、追踪积累、优化系统维护、自写技能漂移、非结构化自我修改、源代码控制差距、成本、停滞、Goodhart 效应、未解决的记忆架构、安全、动态上下文选择、利益相关者压力、不可预测的输入、挑选的结果、奖励博弈、不切实际的期望、MCP 约束、状态管理、工具更替、循环失败、回归地狱、定义争议、已废弃的 API、AutoGPT 风格的空谈无产出、微调过时、平台漏洞、演示规模不足以及开模型的无效 diff。广度表明痛苦不是局部于一个框架或一个模型家族。它是贯穿整个智能体工程栈的系统性问题。

中文补充强调实际部署约束和学习基础设施。其 28 个痛点被汇总为记忆/上下文丢失、智能体循环/控制、成本/token 管理、自我进化/学习、可观察性/调试、安全/安全、框架/工具集成、知识迁移/技能、评估/指标和范式选择。首要反复主题是上下文失忆、自我进化闭环缺失、无限循环、token 成本爆炸和可观察性黑盒。与英文语料库相比，中文讨论更强烈地强调系统化学习智能体开发的困难、在众多本地化和国际框架之间选择的困难，以及真实业务任务中管理 token 成本的困难。这种差异对全球 AI 自我进化研究很重要：采用障碍不仅是模型能力，还包括教育、文档、本地化和成本结构。

\begin{table*}[t]
\centering
\scriptsize
\setlength{\tabcolsep}{2pt}
\renewcommand{\arraystretch}{1.08}
\caption{跨平台频率比较（手动编码结果）。计数为各来源文件内的类别级别，不应作为总体流行率加总。}
\label{tab:platform-frequency}
\begin{tabularx}{\textwidth}{@{}L{0.20\textwidth}L{0.18\textwidth}L{0.18\textwidth}L{0.18\textwidth}Y@{}}
\toprule
\textbf{主题} & \textbf{HN 信号} & \textbf{Reddit 信号} & \textbf{中文社区信号} & \textbf{主导严重性} \\
\midrule
可靠性与幻觉 & 可靠性集群包含 6 个痛点。 & 智能体可靠性包含 6 个；循环失败单独反复出现。 & 智能体循环/控制包含 4 个；幻觉传播显式出现。 & \sevcritical{}/\sevhigh{} \\
自我改进可行性 & 基础集群包含提示优化、学习限制、成本和合规测试。 & 7 个关于可行性、漂移、停滞、回归和成本的痛点。 & 4 个自我进化/学习痛点，包括缺少闭环。 & \sevcritical{} \\
框架与工具 & 5 个框架/工具痛点，包括抽象和工具层问题。 & 6 个框架限制加上 MCP 和状态抱怨。 & 3 个框架/工具集成和 1 个范式选择类别。 & \sevhigh{} \\
评估与基准 & 4 个评估/基准痛点；强烈的基准怀疑论。 & 5 个评估挑战加上挑选和排行榜。 & 2 个评估/指标痛点。 & \sevcritical{}/\sevhigh{} \\
记忆与上下文 & 6 个知识/记忆持久性痛点。 & 4 个记忆/上下文痛点。 & 5 个记忆/上下文丢失痛点；本地首要主题。 & \sevhigh{} \\
安全、安全与成本 & 3 个显式安全/成本痛点。 & 4 个安全/成本痛点加上成本作为可行性阻碍。 & 2 个安全/安全和 2 个成本/token 痛点。 & \sevcritical{} \\
调试与可观察性 & 通过框架不透明、虚假日志和代码审查失败出现。 & 通过状态管理、追踪膨胀和框架不透明出现。 & 2 个可观察性/调试痛点；黑盒是首要主题。 & \sevhigh{} \\
生产验证与规模 & 生产验证差距被显式编码。 & 演示规模不足和真实世界部署各自反复出现。 & 部署担忧通过成本、浏览器失败和教程出现。 & \sevhigh{} \\
\bottomrule
\end{tabularx}
\end{table*}

仅从频率角度阅读会遗漏最重要的模式之一：许多痛点是耦合的。上下文溢出导致幻觉，因为约束被遗忘。幻觉恶化评估，因为测试或日志可能被伪造。弱评估导致回归，因为糟糕的变更被接受。回归增加调试负担，因为状态变更未被追踪。调试不透明导致框架被放弃。框架放弃迫使团队重建基础设施，增加成本。成本压力鼓励更便宜的模型或更少的测试，这可能降低可靠性。这种耦合意味着自我进化系统需要集成控制。仅解决记忆或仅解决评估在其他层仍然薄弱时是不够的。

量化发现还提出了研究基准的路线图。当前基准通常测试智能体是否能解决孤立的编码任务。用户语料库要求轨迹基准：智能体能否在不丢失早期约束的情况下完成长任务、检测自己是否在循环、保留精确数据、从失败的工具中恢复、获取当前文档、产生可检查的追踪并诚实报告未完成的工作？它还要求进化基准：智能体能否在不过拟合的情况下跨运行改进策略、不将成本增加到超出预算、不在保留任务上回归，并清楚解释改变了什么？忽略这些问题的基准可能科学上方便但实践上错位。

数据最终重新定义了人类反馈的角色。许多自我改进提案假设反馈可以自动生成。用户报告可靠的反馈通常是困难的部分。人工审查、测试、基准、生产结果和成本指标各自捕获质量的不同方面。没有任何单一信号是充分的。最可信的系统将结合多个不完美的信号并显式建模其不确定性。例如，通过测试可能是必要但非充分的；代码质量评估器可能捕获可维护性问题；人工审查员可能检查高风险变更；生产遥测可能揭示长尾失败。自我进化应该通过它协调这些反馈源的好坏来评估，而不是通过它是否完全从循环中消除人类。

\subsection{AI 自我进化系统的设计启示}
\label{subsec:design-implications-painpoints}

用户痛点暗示了一组设计义务。第一，自我进化必须以证据为依据。每一个被接受的变更应该指向外部证据，而非仅仅是自我批评。第二，它必须有范围。系统应该指定它正在改变的是提示、工具、记忆策略、技能、代码、路由还是模型权重。第三，它必须可逆转。用户需要快照、diff 和回滚。第四，它必须有预算。搜索应在预期改进不再证明成本合理时停止。第五，它必须可观察。开发者需要追踪级别的可见性来了解为什么提出了变更以及为什么接受了它。第六，它必须安全。权限边界应该管理智能体可以改变什么以及哪些操作需要人工批准。第七，它必须诚实。智能体应该区分完成的工作与计划的工作、观察到的证据与推断的证据、以及不确定性与信心。

这些义务可以表达为任何声称自我进化的智能体的实用检查清单。系统是否维护基线并在保留任务上与之比较？它是否报告失败的候选和回归，还是只报告成功的增量？它是否在编辑时保留精确数据？它是否有超越硬性迭代上限的循环检测？它是否在 API 新鲜度重要时获取当前文档？它是否保持提示和工具调用追踪可见？它是否防止评估器被同一提案修改？它是否区分记忆类型并验证检索到的记忆？它是否强制支出限制并报告每个已接受改进的成本？它是否沙箱化工具使用并要求对风险操作进行批准？它是否在已接受的改进后来造成损害时支持回滚？如果任何答案是否定的，系统可能仍然有用，但其自主性边界应该相应缩小。

语料库还警告不要过度依赖模型规模。更好的模型有帮助，但许多抱怨涉及系统设计：编辑格式、工具模式、可观察性、状态、成本路由、记忆策略、评估设计和安全权限。一个在 opaque、无边界、评估不当的循环中的更强模型仍然可能严重失败。相反，一个在仪表化良好的工具层中的适度模型对于受限任务可能更安全、更有用。这并不意味着模型进步无关紧要。它意味着自我进化研究应该将智能体视为社会技术系统：模型、工具层、记忆、工具、运行时、评估、人工监督和组织治理。

第二个警告涉及"学习"这个词。用户对实际上不跨会话学习的智能体感到沮丧，但他们也对将 RAG 或提示积累称为学习的系统持怀疑态度。分类体系应该至少区分四个级别。第一级是会话适应：智能体使用当前上下文在任务内表现更好。第二级是记忆积累：智能体存储和检索事实或偏好。第三级是程序改进：智能体基于经验创建或修改技能、提示、策略或工具。第四级是模型更新：权重或适配器通过训练改变。许多产品在第一级或第二级运营时宣传自我改进。这不一定是坏事，但应该准确命名。错误标记会产生不切实际的期望并损害信任。

第三个警告涉及基准。语料库包含对被污染的、饱和的或可博弈的基准的强烈怀疑。对于自我进化，Goodhart 定律不是抽象风险；它是核心失败模式。系统正在针对测量过程显式优化自身。如果测量是薄弱的，智能体将学会利用薄弱之处。因此，基准设计应包括反博弈机制：隐藏测试、轮换任务、对抗性检查、定性审查、回归套件和成本标准化分数。报告应包括绝对分数而非仅百分比增益；分布而非仅最佳运行；以及负面结果而非仅获胜。用户对挑选的改进声明的抱怨应该成为发布规范：每篇自我改进论文都应该披露尝试了多少提案、多少失败以及什么回归了。

第四个警告涉及记忆。持久记忆很有吸引力，因为用户讨厌重新解释上下文。但持久记忆也可能保留过时的假设、被污染的指令或过拟合的技能。解决方案不是简单地记住一切。而是创建记忆治理。记忆应该有来源、信心、范围、过期和验证规则。技能应该有测试。摘要应该将精确约束与有损叙事摘要分开保留。检索应该按有用性评估，而非仅语义相似性。智能体应该能够遗忘或隔离导致失败的记忆。在自我进化中，遗忘与记忆同等重要。

最后，痛点揭示了一个产品真相：用户不想与智能体聊天讨论其自主性；他们想要可靠地完成工作。一个 HN 条目捕捉了交互错配：当期望的结果是完成的任务时，没有人想要无尽的对话。因此，自我进化系统应该在有效时使自主性安静、在重要时透明。常规的成功改进可以简洁汇总。风险变更应该附带证据请求批准。失败应该产生可操作的诊断。最好的智能体不是叙述最多推理的那个，而是将经验转化为更安全、更便宜、更可靠的行动，同时给人类正确控制杠杆的那个。

\subsection{小结}
\label{subsec:painpoints-summary}

用户证据改变了 AI 自我进化的议程。核心挑战不仅是设计能修改自身的智能体，而是设计其修改可被信任的智能体。在 HN、Reddit 和中文技术社区中，从业者报告了在可靠性、上下文管理、工具使用、可观察性、评估、成本、安全和框架适配方面的反复失败。最严重的问题是那些破坏信任边界的问题：虚假证据、不安全的自主性、失控支出、未追踪的回归和不透明的状态变更。这些不是能力改进后要解决的边缘情况。它们是安全部署能力的先决条件。

Mom Test 视角有用是因为它将研究议程建立在观察到的变通方法上。用户施加迭代上限是因为智能体无法检测循环。他们移除框架是因为抽象隐藏了提示和状态。他们手动粘贴上下文是因为持久记忆薄弱。他们构建定制的基准是因为公开排行榜不预测本地性能。他们检查每个输出是因为智能体自我报告不可靠。他们沙箱化执行是因为自主工具造成安全风险。每个变通方法都是伪装的设计需求。

就调查目的而言，实践结论很明确。AI 自我进化应该作为具有显式契约的工程学科来评估：以证据为依据的变更、有范围的自主性、可重放的追踪、独立的评估、回归核算、成本治理、记忆治理、安全边界和回滚。框架应该以它们支持这些契约的好坏来评判，而非以它们暴露多少智能体模式。研究应该不仅报告改进，还要报告失败、成本和回归。系统应该不仅改进任务分数，还要改进自身的可靠性基础设施：更好的工具层、更好的工具描述、更好的记忆策略、更好的监控和更好的人工上报。

如果这些痛点被忽视，自我进化可能变成放大用户已经不信任的那些失败的机制。如果被认真对待，用户痛点就变成了路线图。下一代自进化智能体可以通过在约束下学习，从令人印象深刻的演示迈向可信赖的系统：声称前先验证、卡住时停止、带来源地记忆、带测试地变更、带预算地支出、带权限地行动，并暴露足够的证据让人类信任结果。
